\documentclass[10pt,a4paper]{article}
\usepackage[margin=23mm]{geometry}
\usepackage{amsmath,amssymb,amsthm,booktabs,hyperref,fancyhdr,enumitem}
\hypersetup{colorlinks=true,linkcolor=blue,urlcolor=blue}
\pagestyle{fancy}\fancyhf{}\fancyhead[L]{UD: clock--plant compatibility}\fancyhead[R]{CPC v0.1}\fancyfoot[C]{\thepage}
\setlength{\headheight}{14pt}
\newcommand{\norm}[1]{\left\|#1\right\|_\infty}
\newtheorem{theorem}{Theorem}
\title{Clock--Plant Compatibility and Filtered Readout\\for the Proposed UD Diamond Apparatus}
\author{Ben W. Mayes research workstream\\Codex-assisted mathematical development}
\date{20 October 2026 --- Version 0.1}
\begin{document}
\maketitle
\begin{abstract}
We derive an exact compatibility interval linking acquisition period, filter history, actuator rate and the LC plant. Slowing the target cannot be treated as a free remedy for converter noise: under the inherited normalized plant bound, every LC product must grow as the square of the physical model clock. A filtered-readout theorem charges the absolute first moment of its history. A preliminary manufacturer audit motivates an alternative input/command error allocation that still fits the earlier Gauss-discrimination tolerance. All hardware admission remains HOLD.
\end{abstract}
\noindent\textbf{Status.} DERIVED CONDITIONAL; CANDIDATE component route; development, no registry entry; constructor verification only; independent review pending; no experimental execution; physical promotion 0.
\section{The clock and plant must be designed together}
For each canonical pair use $\Phi_{0c}Q_{0c}=S_{\rm action}$ and
$t_0=S_{\rm action}/E_{\rm energy}$. For $L_c,C_c>0$, the normalized free plant has
\begin{equation}
\alpha_c=\frac{\Phi_{0c}^2}{E_{\rm energy}L_c},\quad
\beta_c=\frac{Q_{0c}^2}{E_{\rm energy}C_c},\quad
\lambda=\max_c(\alpha_c,\beta_c).
\end{equation}
The common action has physical units; it is not the structural coefficient-ledger object with a similar symbol.
\begin{theorem}[Exact clock--plant compatibility]
For fixed $t_0$ and physical components,
\begin{equation}
\alpha_c\beta_c=\frac{t_0^2}{L_cC_c},\qquad
\min_{\Phi_{0c}Q_{0c}=S_{\rm action}}\lambda
=\max_c\frac{t_0}{\sqrt{L_cC_c}} .
\end{equation}
Thus $\lambda\leq1$ is achievable by scale selection if and only if
$t_0\leq\min_c\sqrt{L_cC_c}$.
\end{theorem}
\begin{proof}
Multiply (1), then use $\max(\alpha_c,\beta_c)\geq\sqrt{\alpha_c\beta_c}$.
Equality on every channel is attained by
$\Phi_{0c}^2=S_{\rm action}\sqrt{L_c/C_c}$ and
$Q_{0c}^2=S_{\rm action}\sqrt{C_c/L_c}$.
These choices preserve the same common action across channels.
\end{proof}
With physical output period $T_p$, absolute filter-history moment $\mu_p$, additional electronics latency $\ell_p$, and actuator time constant $\tau_p$, the SAC allocation requires
\begin{equation}
400000\max(T_p,\mu_p+\ell_p,\tau_p)
\leq t_0\leq\min_c\sqrt{L_cC_c}.
\end{equation}
The lower bound is a sufficient timing allocation once these quantities are justified, not a measured capability. A nonempty interval is necessary for this particular contract. Common-action calibration and actual loading remain separate gates.
\newpage
\section{A filtered ADC is a history measurement}
\begin{theorem}[Filtered-readout error]
Suppose physical state speed is at most $M/t_0$. Let
$y(t)=\sum_jw_j[z(t-s_j)+e_j]$, where $s_j\geq0$,
$\sum_jw_j=1$, and $\norm{e_j}\leq\nu_{\rm raw}$. Then
\begin{equation}
\norm{y(t)-z(t)}
\leq\frac{M}{t_0}\sum_j|w_j|s_j+
\nu_{\rm raw}\sum_j|w_j|.
\end{equation}
\end{theorem}
\begin{proof}
Subtract $z(t)$, use the unit sum, bound each state difference by $Ms_j/t_0$, and apply the triangle inequality.
\end{proof}
For nonnegative weights the absolute first moment is the mean age and the noise-bound amplification is one. For signed FIR weights a signed group delay can cancel. Weights $(2,-2,1)$ at ages $(0,1,2)$ have unit sum and zero signed mean age, but absolute moment four.

Charge output holding, electronics latency and filter history once as
$M(T_p+\ell_p+\mu_p)/t_0$ in the command-age budget; charge static/noise reconstruction separately. Published group delay can serve as an absolute moment only with a justified nonnegative, unit-gain history and correct aperture/availability timing. Analog filtering, startup tails, crosstalk and unrecorded timestamps need bounds. Typical RMS noise is not a simultaneous hard bound.
\subsection*{Arbitrary slowing also costs control accuracy}
Let $r=\min_c\sqrt{L_cC_c}$ and $A=T_p+\ell_p+\mu_p+\tau_p$. Within the existing conservative SAC formula $K_U=L_F+\lambda$, use the smallest possible scalar plant bound $\lambda=t_0/r$ from (2). Its steady tracking-budget expression is
\begin{equation}
B(t_0)=\frac{L_FMA}{t_0}+\frac{MA}{r}
+L_F\nu+\frac{\nu t_0}{r}+q+\epsilon_a,\quad L_F=97/4 .
\end{equation}
For $\nu>0$, arithmetic--geometric mean gives
\begin{equation}
B(t_0)\geq\frac{MA}{r}+L_F\nu+q+\epsilon_a+
2\sqrt{\frac{L_FMA\nu}{r}},\quad
t_0^*=\sqrt{\frac{L_FMAr}{\nu}} .
\end{equation}
Other timing, range and plant restrictions can exclude this unconstrained optimum. Bounds must be revalidated after any change of physical scales. Even dropping the positive terms, $A\geq1.6$ ms and $r\leq1$ s give $B\geq81/6250=0.01296$, above the inherited tracking allocation $1/400$ for every clock choice.

\textbf{Scope.} $B$ is the value of this conservative sufficient certificate, not a lower bound on actual apparatus error. Sharper sensitivity estimates, predictors, different readouts or different plants can change the conclusion. Equation (3) under $\lambda\leq1$ is an exact compatibility restriction, not a universal no-go for LC simulation.
\newpage
\section{Manufacturer audit and an alternative budget}
These are supplier facts under stated conditions, not apparatus observations. Documents were accessed on 20 October 2026.
\begin{center}
\begin{tabular}{p{0.23\linewidth}p{0.66\linewidth}}
\toprule
Candidate & Scoped finding and implication\\
\midrule
AD5791B DAC [1] & 20-bit output; maximum INL one LSB for $\pm10$ V references at $0$--$105^\circ$C. Settling depends on step size and buffer; it does not identify a first-order plant.\\
LTC2500-32 ADC [2] & Filtered output is 32 bits; the no-latency composite contains 24 differential bits. Maximum INL is 2 ppm. Averaging at 1 MHz and DF=1024 gives spacing 1.024 ms and listed group delay 512 $\mu$s.\\
ADS127L21 [3] & Alternative filtered ADC route; operating mode, noise and history require a separate audit before selection.\\
\bottomrule
\end{tabular}
\end{center}
These are screened candidates, not a purchasing recommendation or an admitted bill of materials. References, buffers, voltage/current drivers, synchronized channel acquisition and the canonical plant remain unspecified.

For nominal DAC span $[-10,10]$, nearest-code quantization plus one-LSB INL has conservative interior-range allowance
\begin{equation}
\frac{10}{2^{20}}+\frac{20}{2^{20}}
=\frac{15}{524288}<\frac1{25000}.
\end{equation}
Endpoint representability also needs checking; SAC commands lie well inside this span. Gain, offset, reference, noise, drift, computation and driver errors must fit the remainder after calibration. No calibration is supplied.

A 2-ppm-of-span ADC INL on normalized $[-1,1]$ contributes $4/10^6$. Consider the \emph{alternative} supplied bounds $\nu=1/100000$ and $q=1/25000$, retaining the other SAC values. Its tracking and total residual are
\begin{equation}
b=\frac{29383}{16000000},\qquad
p+\max(\sigma,b)=\frac{149383}{16000000}
=0.0093364375<0.01 .
\end{equation}
Slack is $10617/16000000$. The prior Gauss theorem retains conservative $\delta=1/100$ and its preparation, endpoint reconstruction and clock-error requirements. With the same hold radius and ten-time-constant preload, release tracking is below $6449/8000000<1/1000$. This alternative does not correct or supersede SAC's earlier valid example.
\subsection*{A screen that prevents premature component selection}
The illustrated averaging output spacing alone requires
$t_0\geq400000(1.024\text{ ms})=2048/5$ s. Under $\lambda\leq1$, every pair needs
\begin{equation}
L_cC_c\geq4194304/25\ {\rm s}^2 .
\end{equation}
The illustrative choice $t_0=512$ s requires $L_cC_c\geq262144$ s$^2$ and makes the $1/100$ prefix last 5.12 s. This does not establish feasible components. If each pair has $L_cC_c\leq1$ s$^2$, the clock interval is empty.
\newpage
\section{An instrumented identification plan}
The next physical return is a calibration packet, not an execution claim.
\begin{enumerate}[leftmargin=*,itemsep=4pt]
\item \textbf{Clock and plant.} Return actual $L_c,C_c$, loading, action products and energy scales with units and joint intervals. Evaluate (3) on worst-case intervals. Reject a route with no admissible clock.
\item \textbf{Readout history.} Freeze mode, filter coefficients, aperture times and output timestamps. Return the absolute first moment, coefficient absolute sum, startup history, channel skew and total acquisition-to-command latency. Unknown history yields HOLD.
\item \textbf{Static accuracy.} Measure traceable code/voltage maps, gain, offset, drift and crosstalk. Supply a finite cover and separately justified regularity bounds. Add simultaneous noise/reconstruction uncertainty at required joint coverage; do not infer peak bounds from RMS statistics.
\item \textbf{Actuator model.} Independently measure actual commands $c$, outputs $a$ and physical derivatives. Over each held-command segment bound
\[
r_a(t)=\dot a(t)/\omega_p+a(t)-c(t).
\]
SAC requires $\norm{r_a}\leq\epsilon_a$. Charge derivative reconstruction, output/command uncertainty, rate uncertainty and regularity times cover radius. Include one-sided endpoints, code-change glitches, tails, cross-channel loads, range and slew.
\item \textbf{Release and blinded test.} Bound hold radius, prehistory speed and release tracking; reset the clock at release and record holding work. After premises and non-constructor review are on record, evaluate nominal and fault runs under the frozen rule.
\end{enumerate}
No component was purchased, researcher contacted, or experiment executed. Exact calculations and inherited-budget replay verify algebra only; hardware remains HOLD. Non-constructor review and physical promotion have not occurred.
\subsection*{Source and dependency record}
\begin{enumerate}[leftmargin=*,itemsep=3pt]
\item Analog Devices, \emph{AD5791}, Rev. F, Table 3 p.3; terminology p.17.
\href{https://www.analog.com/media/en/technical-documentation/data-sheets/AD5791.pdf}{Manufacturer data sheet (analog.com).}
\item Analog Devices/Linear Technology, \emph{LTC2500-32}, Rev. B, electrical characteristics p.4 and Table 2 p.33.
\href{https://www.analog.com/media/en/technical-documentation/data-sheets/250032fb.pdf}{Manufacturer data sheet (analog.com).}
\item Texas Instruments, \emph{ADS127L21}, Rev. C, August 2025.
\href{https://www.ti.com/lit/ds/symlink/ads127l21.pdf}{Manufacturer data sheet (ti.com).}
\end{enumerate}
Internal dependencies: CAM/CD-1, FCAL and SAC v0.1 retain their assumptions and statuses. Citation does not strengthen physical evidence. The elementary clock-product and filter inequalities carry no general novelty claim.
\vfill\noindent\textbf{Claim footer.} DERIVED CONDITIONAL; CANDIDATE component route; hardware HOLD; development, no registry entry; constructor checks only; independent review pending; execution NONE; physical promotion 0.
\end{document}
